\documentclass[aspectratio=1610]{beamer}

\usetheme{osx}
\btosxwebsite[didierverna.info]{https://didierverna.info}
\btosxx{didierverna}
\btosxfacebook{didier.verna}
\btosxlinkedin{didierverna}
\btosxcopyright[string=all rights reserved.]
\btosxbuild{Build 1.0}
\setlength\BTOSXFolderSep{30pt}
\definecolor{BTOSXCurrentSectionColor}{named}{yellow}


\title{Beamer Theme OSX}
\subtitle{An OSX 10.5 (Leopard) look for your slides}
\author{Didier Verna}
\institute{\em Resistance is futile.}
\date{\today}


\begin{document}

\frame{\titlepage}

\begin{frame}{Outline}
  \tableofcontents
\end{frame}

\section[First]{First Section}
\subsection[First First]{First Subsection of the First Section}
\begin{frame}{A Frame}
  Some text bla bla bla.
  \begin{block}{A block}
    With some block contents
  \end{block}
\end{frame}

\subsection[Second First]{Second Subsection of the First Section}
\begin{frame}{Another Frame}
  Some text bla bla bla.
  \begin{alertblock}{An alterted block}
    With some block contents
  \end{alertblock}
\end{frame}

\section[Second]{Second Section}
\subsection[First Second]{First Subsection of the Second Section}
\begin{frame}[t]{Geometry values}
  hoffset = \the\hoffset, voffset = \the\voffset\\
  oddsidemargin = \the\oddsidemargin\\
  topmargin = \the\topmargin\\
  headheight = \the\headheight, headsep = \the\headsep\\
  textheight = \the\textheight, textwidth = \the\textwidth\\
  marginparsep = \the\marginparsep, marginparwidth = \the\marginparwidth\\
  footskip = \the\footskip\\

  marginparpush = \the\marginparpush\\
  paperwidth = \the\paperwidth, paperheight = \the\paperheight\\

  beamer@leftsidebar = \the\csname beamer@leftsidebar\endcsname\\
  beamer@rightsidebar = \the\csname beamer@rightsidebar\endcsname\\
  beamer@leftmargin = \the\csname beamer@leftmargin\endcsname\\
  beamer@rightmargin = \the\csname beamer@rightmargin\endcsname
\end{frame}

\subsection[Second Second]{Second Subsection of the Second Section}
\begin{frame}[t]{Margin test}
  And this is some very long text in order to test for the contents margins,
  so that, you know, it all remains within the bounds of this cute little Mac
  OSX like window that I've come up with in this super cool Beamer theme
  which, BTW, requires a lot of wizardry and a fair amount of TikZ mastery,
  because I gotta tellya my friend, all of this is quite tricky, with a 1000
  pages or so long manual, and such an amount of\ldots syntax! On the other
  hand, I designed this theme 10 years ago, and if I had to do it all over
  again today, I probably would just ask Claude to do it for me.
\end{frame}

\end{document}

%%% Local Variables:
%%% mode: latex
%%% TeX-master: t
%%% End:
